\chapter{Tổng quan về đề tài}
 Nội dung chương 1trình bày sơ lược về đề tài, bối cảnh hiện nay và làm rõ lý do đề tài được
 chọn. Từ đó, xác định mục tiêu cần đạt được khi thực hiện và phạm vi của đề tài.
\section{Lý do chọn đề tài}
Trong bối cảnh nền kinh tế số toàn cầu đang bước vào giai đoạn phát triển bùng nổ, việc lựa chọn nghiên cứu đề tài Hệ thống gợi ý sản phẩm trên các sàn thương mại điện tử xuất phát từ sự hội tụ của ba động lực chính: tính cấp thiết của thị trường, giá trị kinh tế khổng lồ và bước ngoặt công nghệ mang tính cách mạng. Chúng ta đang sống trong một kỷ nguyên mà khả năng sản xuất và cung ứng hàng hóa đã vượt xa khả năng xử lý thông tin của con người. Với doanh số thương mại điện tử toàn cầu ước tính đạt 6,33 nghìn tỷ USD vào năm 2024 và dự kiến tiếp tục tăng trưởng lên 6,86 nghìn tỷ USD vào năm 2025, người tiêu dùng đang bị bủa vây bởi một biển hàng hóa. Đặc biệt, khi thương mại di động chiếm tới 57\% tổng doanh số và dự báo đạt 2,51 nghìn tỷ USD vào năm 2025 nhưng không gian hiển thị trên màn hình điện thoại vô cùng khan hiếm~\cite{1}. Sự mâu thuẫn giữa số lượng sản phẩm khổng lồ và giới hạn nhận thức của người dùng đã tạo ra nhu cầu bức thiết cho những cơ chế lọc thông minh, có khả năng thấu hiểu và dự đoán nhu cầu trước khi người dùng nhận thức được chúng.\\

Về mặt kinh tế, hệ thống gợi ý đã chuyển mình từ một tính năng phụ trợ trở thành động cơ doanh thu quyết định sự sống còn của doanh nghiệp. Các dữ liệu thực nghiệm từ những tập đoàn công nghệ hàng đầu đã minh chứng rõ nét cho luận điểm này. Amazon, gã khổng lồ bán lẻ trực tuyến, tạo ra tới 35\% tổng doanh thu từ các cơ chế gợi ý cá nhân hóa~\cite{2}. Trong lĩnh vực nội dung số, Netflix tiết kiệm khoảng 1 tỷ USD mỗi năm và giữ chân người dùng nhờ vào hệ thống gợi ý đóng góp tới 75-80\% thời lượng xem~\cite{2}. Đối với các doanh nghiệp vừa và nhỏ, việc áp dụng cá nhân hóa bằng AI có thể giúp tăng doanh thu từ 10-40\% và cải thiện tỷ lệ chuyển đổi trung bình khoảng 26\%~\cite{3}. Ngược lại, sự thiếu vắng một hệ thống gợi ý hiệu quả dẫn đến cái giá rất đắt: hiện tượng bỏ dở tìm kiếm xảy ra ngày càng nhiều, khi người dùng rời bỏ trang web vì không tìm thấy sản phẩm ưng ý, đang gây thiệt hại cho ngành bán lẻ toàn cầu hơn 2.000 tỷ USD mỗi năm~\cite{4}. Do đó, nghiên cứu xây dựng hệ thống gợi ý không chỉ là bài toán kỹ thuật mà là giải pháp tối ưu hóa dòng tiền và lợi nhuận thiết yếu.\\

Xét trên khía cạnh tâm lý học hành vi, đề tài này giải quyết vấn đề cốt lõi của người tiêu dùng hiện đại: sự tê liệt khi lựa chọn. Các nghiên cứu chỉ ra rằng khi đối mặt với quá nhiều tùy chọn, não bộ con người rơi vào trạng thái mệt mỏi nhận thức, dẫn đến việc trì hoãn hoặc từ bỏ mua sắm. Hơn nữa, hành vi người dùng đang chuyển dịch mạnh mẽ từ tìm kiếm dựa trên từ khóa sang khám phá dựa trên gợi ý, đặc biệt là tại các thị trường mới nổi như Việt Nam. Với quy mô thị trường thương mại điện tử Việt Nam, doanh thu của bốn nền tảng thương mại điện tử Shopee, TikTok Shop, Lazada và Tiki trong quý I 2024 là 79,12 nghìn tỷ đồng, trong đó TikTok Shop chiếm lĩnh 23,2\% thị phần chỉ sau thời gian ngắn\cite{17}, người tiêu dùng ngày càng mong đợi các trải nghiệm mua sắm mang tính giải trí, ngẫu hứng và được cá nhân hóa cao độ. Hệ thống gợi ý chính là chìa khóa để đáp ứng kỳ vọng này, đóng vai trò như một người trợ lý ảo thấu hiểu tâm lý, giúp điều hướng người dùng trong mê cung hàng hóa và kích thích hành vi mua sắm.\\

Cuối cùng, lý do quan trọng để thực hiện đề tài này ngay tại thời điểm hiện tại nằm ở sự tiến hóa vượt bậc của công nghệ. Lĩnh vực hệ thống gợi ý đang trải qua một cuộc cách mạng với sự xuất hiện của Generative AI và các Mô hình ngôn ngữ lớn. Các phương pháp truyền thống như Lọc cộng tác đang dần bộc lộ hạn chế trong việc xử lý vấn đề khởi động lạnh và thiếu khả năng giải thích lý do gợi ý. Sự tích hợp LLM vào hệ thống gợi ý đang mở ra hướng đi mới, cho phép hệ thống không chỉ phân tích hành vi click chuột mà còn đọc hiểu được ý định của người dùng thông qua ngôn ngữ tự nhiên, giúp cải thiện các chỉ số dự đoán quan trọng như AUC lên tới 5\% và NDCG lên tới 170\% trong các tác vụ gợi ý tuần tự~\cite{5}. Hơn thế nữa, sự phát triển của các AI Agents trong năm 2025 hứa hẹn đưa cá nhân hóa lên một tầm cao mới, nơi AI có thể tự chủ lập kế hoạch và thực hiện các tác vụ mua sắm phức tạp thay cho con người.\\

Hệ thống gợi ý sản phẩm trong sàn thương mại điện tử đang là một lựa chọn mang tính chiến lược, đón đầu xu hướng công nghệ của giai đoạn 2025-2026. Nó không chỉ có ý nghĩa khoa học trong việc ứng dụng các thành tựu AI mới nhất mà còn mang lại giá trị thực tiễn to lớn, giúp giải quyết bài toán hiệu quả kinh doanh cho doanh nghiệp và nâng cao trải nghiệm cho người tiêu dùng trong nền kinh tế số.

\section{Mục tiêu đề tài}
Từ những lợi ích thiết thực trên, nhóm đã lên kế hoạch thực hiện một hệ thống gợi ý sản phẩm thông minh thế hệ mới cho thương mại điện tử, hoạt động như một trợ lý mua sắm ảo có khả năng thấu hiểu đa chiều. Hệ thống nhằm giải quyết bài toán quá tải lựa chọn của người tiêu dùng, giúp người dùng tìm được sản phẩm ưng ý nhanh nhất với ít thao tác nhất.\\

Để đạt được mục tiêu tổng quát nêu trên, nhóm nghiên cứu tập trung vào việc xây dựng khả năng thấu hiểu dữ liệu đa phương thức, khắc phục hạn chế của các hệ thống truyền thống vốn chỉ dựa trên từ khóa hoặc định danh sản phẩm. Hệ thống được thiết kế để xử lý và trích xuất đặc trưng sâu từ cả hình ảnh và văn bản mô tả chi tiết, cho phép nhận diện chính xác các yếu tố về kiểu dáng, màu sắc và phong cách để đưa ra các gợi ý có độ tương đồng cao về mặt thị giác và ngữ nghĩa. Bên cạnh đó, đề tài cũng chú trọng thiết lập cơ chế theo dõi và phân tích chuỗi tương tác theo thời gian, giúp hệ thống nắm bắt được sở thích nhất thời và sự thay đổi trong gu thẩm mỹ của người dùng, đảm bảo các sản phẩm được đề xuất luôn phù hợp với ý định mua sắm tại thời điểm hiện tại.

Trong quá trình xác định giải pháp công nghệ cho hệ thống, nhóm nghiên cứu đã tiến hành phân tích và so sánh kỹ lưỡng giữa các hướng tiếp cận phổ biến hiện nay nhằm tìm ra phương án tối ưu nhất. Đối với các phương pháp truyền thống như Lọc cộng tác hay Lọc theo nội dung, ưu điểm nằm ở sự đơn giản và chi phí tính toán thấp, tuy nhiên chúng bộc lộ nhược điểm chí mạng về vấn đề khởi động lạnh cũng như khả năng xử lý dữ liệu thưa kém, dẫn đến việc không thể đưa ra gợi ý chính xác cho người dùng mới hoặc sản phẩm mới chưa có lịch sử tương tác. Trong khi đó, các mô hình học sâu thuần túy tuy cải thiện được độ chính xác nhờ khả năng nắm bắt các mối quan hệ phi tuyến tính trong chuỗi hành vi, nhưng lại thường hạn chế trong việc hiểu sâu ngữ nghĩa đa phương thức của hình ảnh và văn bản, đồng thời hoạt động như một hộp đen thiếu tính minh bạch. Riêng đối với các Mô hình ngôn ngữ lớn, mặc dù sở hữu năng lực thấu hiểu ngữ cảnh và suy luận ý định vượt trội, nhưng việc triển khai độc lập cho toàn bộ quy trình lại vấp phải rào cản lớn về chi phí tính toán khổng lồ và độ trễ cao, khiến chúng trở nên kém khả thi khi phải xử lý xếp hạng hàng triệu sản phẩm trong thời gian thực. Chính vì những lý do trên, nhóm quyết định lựa chọn giải pháp kiến trúc lai hiện đại, thiết lập một đường ống xử lý phân tầng chặt chẽ. Hệ thống sử dụng mô hình Two-Tower làm bộ lọc sơ cấp để tối ưu hóa tốc độ truy xuất ứng viên quy mô lớn, kết hợp với mô hình Cross-Encoder ở tầng trung gian để đánh giá mức độ giao thoa đặc trưng chi tiết, và cuối cùng tận dụng khả năng thấu hiểu sâu sắc của LLM cùng kỹ thuật RAG ở tầng tinh chỉnh tối hậu. Lựa chọn này được đánh giá là ưu việt nhất vì nó đạt được sự cân bằng lý tưởng: vừa đảm bảo hiệu năng vận hành nhanh chóng của hệ thống lớn, vừa tận dụng được trí tuệ của AI để giải quyết bài toán khởi động lạnh và cá nhân hóa trải nghiệm người dùng một cách chính xác nhất.

\section{Đối tượng và Phạm vi nghiên cứu}

Đối tượng nghiên cứu trung tâm của đề tài là các hệ thống gợi ý thế hệ mới trong thương mại điện tử, phản ánh xu hướng chuyển dịch mạnh mẽ từ các phương pháp lọc cộng tác dựa trên định danh đơn thuần sang các mô hình học sâu có khả năng nhận thức ngữ cảnh. Cụ thể, đề tài đi sâu vào phân tích ba khía cạnh cốt lõi cấu thành nên hệ thống. Đầu tiên là việc mô hình hóa chuỗi hành vi người dùng, tập trung khai thác sự diễn biến động của sở thích theo trục thời gian thông qua các trình tự tương tác liên tục như xem, mua hay đánh giá nhằm nắm bắt quy luật chuyển đổi nhu cầu ngắn hạn và dài hạn. Song song với đó là nghiên cứu cơ chế biểu diễn tri thức sản phẩm đa phương thức, cho phép máy tính thực sự thấu hiểu nội dung hàng hóa nhờ việc hợp nhất các luồng thông tin phức tạp từ thị giác, ngôn ngữ đến dữ liệu cấu trúc thay vì chỉ coi sản phẩm là các mã định danh vô nghĩa. Cuối cùng, đề tài tập trung vào cơ chế suy luận và xếp hạng lại, trong đó ứng dụng các Mô hình ngôn ngữ lớn như những tác nhân thông minh để đánh giá lại độ phù hợp dựa trên ngữ cảnh truy vấn phức tạp, qua đó đảm bảo kết quả gợi ý đạt được tính chính xác cao đi kèm với sự minh bạch trong lý do đề xuất.

Phạm vi thực hiện của đề tài được xác định rõ ràng thông qua ba khía cạnh chính bao gồm không gian ứng dụng, dữ liệu, công nghệ. Không gian nghiên cứu tập trung vào bối cảnh các nền tảng Thương mại điện tử, nơi dữ liệu thường mang đặc thù thưa thớt và chứa nhiều nhiễu, với mục tiêu tối ưu hóa trải nghiệm khám phá cho người dùng cuối. Để thực hiện điều này, đề tài giới hạn phạm vi dữ liệu vào việc khai thác nguồn thông tin đa phương thức phong phú, kết hợp chặt chẽ giữa dữ liệu phi cấu trúc như văn bản mô tả, hình ảnh sản phẩm và dữ liệu có cấu trúc gồm lịch sử tương tác gắn nhãn thời gian. Trên nền tảng đó, nghiên cứu ứng dụng các công nghệ tiên tiến của Trí tuệ nhân tạo, trọng tâm là kỹ thuật Học sâu để trích xuất đặc trưng và các Mô hình ngôn ngữ lớn giúp thấu hiểu ý định người dùng.

\section{Định nghĩa bài toán}

Bài toán cụ thể được đặt ra trong nghiên cứu này là bài toán dự đoán và xếp hạng tuần tự dựa trên dữ liệu đa phương thức, nơi hệ thống phải đối mặt với thách thức xác định chính xác nhu cầu tiếp theo của người dùng trong không gian tìm kiếm khổng lồ của thương mại điện tử. Cụ thể, với đầu vào là chuỗi lịch sử tương tác tuần tự theo thời gian kết hợp cùng dữ liệu nội dung đa phương thức phong phú của sản phẩm như hình ảnh và văn bản mô tả, nhiệm vụ đặt ra là xây dựng một hàm mục tiêu tối ưu có khả năng xếp hạng hàng triệu ứng viên sao cho các sản phẩm ở vị trí đầu tiên có xác suất tương thích lớn nhất với ý định hiện tại của khách hàng. Điểm mấu chốt của bài toán không chỉ dừng lại ở việc khớp lệnh hành vi quá khứ, mà còn phải giải quyết triệt để mâu thuẫn giữa tính thưa thớt của dữ liệu tương tác và sự phong phú của dữ liệu ngữ nghĩa, đặc biệt trong các kịch bản khởi động lạnh khi thông tin hành vi hoàn toàn vắng bóng. Do đó, bài toán yêu cầu một cơ chế suy luận phức hợp, có khả năng chuyển đổi các tín hiệu phi cấu trúc thành các biểu diễn vector đồng nhất trong không gian tiềm ẩn để lấp đầy khoảng trống thông tin mà các phương pháp dựa trên ID truyền thống không thể xử lý được.

Quy trình vận hành tổng thể của hệ thống được thiết kế theo kiến trúc đường ống xử lý dữ liệu phân tầng, khởi tạo từ việc tiếp nhận đầu vào phức hợp bao gồm hồ sơ người dùng, chuỗi lịch sử tương tác thời gian thực và dữ liệu đa phương thức của sản phẩm như hình ảnh, mô tả văn bản chi tiết. Luồng thông tin này trước tiên được đưa qua tầng mã hóa đặc trưng để chuyển đổi thành các vector biểu diễn trong không gian tiềm ẩn, phục vụ cho giai đoạn truy xuất ứng viên. Tại đây, mô hình Two-Tower đóng vai trò bộ lọc sơ cấp, thực hiện quét nhanh trên không gian dữ liệu quy mô lớn để sàng lọc ra một tập hợp nhỏ các sản phẩm tiềm năng nhất dựa trên độ tương đồng vector. Ngay sau đó, tập ứng viên này được đưa qua mô hình Bộ mã hóa chéo để thực hiện bước đánh giá luân phiên và tính toán chi tiết mức độ giao thoa đặc trưng giữa người dùng và từng sản phẩm. Danh sách sau khi được chấm điểm và xếp hạng lại một lần nữa sẽ được chuyển giao cho phân hệ Mô hình ngôn ngữ lớn. Tại đây, công nghệ trí tuệ nhân tạo tạo sinh không chỉ đóng vai trò tinh chỉnh thứ hạng lần cuối dựa trên tri thức tổng quát, mà còn đảm nhiệm việc sinh ra các lời giải thích ngữ nghĩa minh bạch cho từng gợi ý. Kết quả đầu ra cuối cùng của hệ thống là một danh sách xếp hạng các sản phẩm có độ phù hợp cao nhất với ý định người dùng, được hiển thị trực quan đi kèm với các thông tin giải thích ngữ nghĩa minh bạch, giúp tối ưu hóa quyết định mua hàng một cách chính xác và thuyết phục.
